\subsection{A nonproduct mixed class detected by the ordinary polarization index}
\label{sec:DE-A3-mixed-nonproduct}

\begin{lemma}[The two integral polarization indices in mixed dimension two]
\label{lem:DE-A3-mixed-polarization-blocks}
\status{Proved; root and monodromy full independent proof/source/code audits passed.}
Let $A$ be an abelian surface of $2$-rank one over an algebraically
closed field $k$ of characteristic two, and let $M$ be ample.
Put $D=\chi(M)$ and define
\[
 a(M)=\operatorname{length}_{\mathbf Z_2}
       \coker\bigl(T_2(\lambda_M):T_2(A)\longrightarrow T_2(A^\vee)\bigr).
\]
These Tate modules both have rank one.  There is an integer $b(M)\ge0$
such that
\[
                         v_2(D)=a(M)+b(M).
\]
The elementary divisors of the integral polarization pairing on the
ordinary Dieudonne summand are $(a(M),a(M))$, and those on its
height-two local-local summand are $(b(M),b(M))$.  In particular,
$\ker\lambda_M$ is killed by $D$.

There is an ample line bundle $N$ on $A^\vee$ satisfying
\[
 \lambda_M^*N\simeq M^D,\qquad
 \lambda_N=D\lambda_M^{-1},\qquad
 \chi(N)=D,\qquad
             a(N)=b(M),\quad b(N)=a(M).
\]
The inverse here is first taken in the rational homomorphism space;
the assertion includes the integrality of the displayed polarization.
\end{lemma}
\begin{proof}
We explain both the finite-flat splitting and the integral pairing.
Over a perfect field the connected-etale sequence of a finite
commutative group scheme splits canonically and functorially:
the reduced subgroup is the etale factor.  There are no nonzero
homomorphisms between a local group and an etale group.  These
statements, including the functoriality needed for transition maps,
are \cite[Lemma~4.44 and Proposition~4.45, p.~67]
{EdixhovenVanDerGeerMoonenAV}.
Apply them at every level of $A[2^\infty]$, and then apply the same
splitting to the Cartier dual of the connected factor.  The compatible
splittings give an actual product of $2$-divisible groups
\[
 A[2^\infty]=G_{\et}\times G_{\rm mult}\times G_{\rm ll},
 \qquad
 G_{\et}\simeq\mathbf Q_2/\mathbf Z_2,\quad
 G_{\rm mult}\simeq\mu_{2^\infty}.
\]
Here $G_{\rm ll}$ and its Cartier dual are connected.  It has height
two and dimension one.  Indeed the etale factor has height one by
the $2$-rank hypothesis; duality and isogeny invariance give the same
height for the multiplicative factor, leaving height two out of the
total height four.  The same reasoning applies to $A^\vee$.
Maps between distinct factors vanish.  For the multiplicative versus
local-local assertion, apply Cartier duality and the local-etale
vanishing.  Thus every homomorphism under consideration preserves
these factors, before any passage to isogeny categories.

Use the covariant Dieudonne lattice, equivalently the $W(k)$-linear
dual of Oda's contravariant lattice.  It is free, and Cartier duality
is its integral $W(k)$-linear dual, with the usual exchanged
Frobenius and Verschiebung actions.  Exactness and the length/order
dictionary are \cite[Theorem~3.14(ii)--(iii) and
Corollaries~3.15--3.16, pp.~87--88; Proposition~3.18, p.~89;
Proposition~3.22, p.~99]{Oda1969Dieudonne}.
The polarization gives a skew-symmetric integral bilinear Riemann
form on this lattice by
\cite[Propositions~1.10 and 1.12, pp.~75--76, and
Proposition~3.24(ii), pp.~100--101]{Oda1969Dieudonne}.
Since $W(k)$ has no $2$-torsion, skew symmetry implies that the
diagonal values are zero.  The form is therefore alternating even
in the present residue characteristic.  The ordinary and local-local
summands are orthogonal, by the vanishing of the cross-factor
homomorphisms just proved.

The ordinary summand pairs its rank-one etale and multiplicative
parts.  Symmetry of the polarization and integral Cartier duality
make their two indices equal; the etale one is precisely $a(M)$.
On the free rank-two local-local lattice, an alternating matrix has
the form
\[
                         \begin{pmatrix}0&c\\-c&0\end{pmatrix},
                         \qquad 0\ne c\in W(k).
\]
Set $b(M)=v_2(c)$.  Its two elementary divisors are both $b(M)$;
this integer is independent of the chosen basis.
The kernel of an isogeny has Dieudonne length equal to the valuation
of its order.  Since $\deg\lambda_M=D^2$, the two blocks give
$2v_2(D)=2a(M)+2b(M)$.  Their finite cokernels are killed by
$2^{\max\{a(M),b(M)\}}$, hence by the two-primary part of $D$.
Exact Dieudonne theory transfers this assertion to the finite
group-scheme kernel.

For an odd prime $\ell$, the kernel is a constant finite abelian
$\ell$-group with a perfect alternating commutator pairing.
If an element has maximal order $\ell^e$, perfection supplies a
partner with pairing of that same order.  The two generate a group
of order $\ell^{2e}$, by pairing any proposed relation with the two
generators.  Thus $e\le v_\ell(D)$.
Together with the two-primary calculation this proves
$D\ker\lambda_M=0$.

Consequently the commutator pairing of $M^D$ is trivial on
$\ker\lambda_M$.  The group-scheme descent criterion
\cite[Corollary~8.11 and its proof, p.~117]
{EdixhovenVanDerGeerMoonenAV}, which allows a nonreduced isotropic
subgroup, gives $N$ with $\lambda_M^*N\simeq M^D$.
It is ample because its pullback by the finite surjection is ample.
Polarizations of pullbacks give
\[
                   \lambda_M\lambda_N\lambda_M=D\lambda_M,
\]
hence the asserted integral inverse formula.
Intersection degrees give
$\chi(N)=\chi(M^D)/\deg\lambda_M=D^3/D^2=D$.
On each of the two blocks, multiplying the inverse matrix by $D$
replaces its elementary-divisor valuation by $v_2(D)$ minus that
valuation.  The two indices of $N$ are therefore $b(M)$ and $a(M)$,
as stated.
\end{proof}

\begin{lemma}[An etale principal source exactly at ordinary index zero]
\label{lem:DE-A3-mixed-index-zero-source}
\status{Proved; root and monodromy full independent proof/source/code audits passed.}
Under the hypotheses of the preceding lemma, the condition $a(M)=0$
is equivalent to the existence of a principally polarized surface
$(B,\Theta_B)$ and an etale isogeny $\phi:B\to A$ of degree $D$
such that
\[
                \phi^*M\equiv D\Theta_B,\qquad
                \operatorname{Nm}_\phi(\Theta_B)\equiv M.
\]
The surface $B$ has $2$-rank one.  No parity restriction is imposed
on $D$, and no product decomposition of $A$ is assumed.
\end{lemma}
\begin{proof}
Suppose first that $a(M)=0$.  Construct $N$ on $A^\vee$ by
Lemma~\ref{lem:DE-A3-mixed-polarization-blocks}.
Its local-local block has index zero, so the two-primary kernel of
$\lambda_N$ is entirely in the ordinary factor.  It is the product
of a multiplicative group and its etale Cartier dual, each of order
$2^{b(M)}$.  Its commutator pairing is trivial on either factor and
perfect between them: restriction to a single factor is a map from
a local multiplicative group to an etale dual, or its reverse, and
must vanish.  Thus its entire multiplicative factor is maximal
isotropic.

Choose this factor in the two-primary kernel and a maximal isotropic
subgroup in each odd-primary kernel.  Their product $H$ is isotropic
of order $D$, and $H^D$ is etale.  The descent theorem gives
\[
       \psi:A^\vee\longrightarrow C=A^\vee/H,\qquad
       \psi^*\Theta_C\simeq N
\]
with $\Theta_C$ principal
\cite[Proposition~8.19 and proof, p.~119]
{EdixhovenVanDerGeerMoonenAV}.
Put $B=C^\vee$, take its dual principal polarization, and set
$\phi=\psi^\vee$.  Isogeny duality gives
$\ker\phi=H^D$ \cite[Theorem~7.5, p.~100]
{EdixhovenVanDerGeerMoonenAV}, so $\phi$ is etale.
The typed polarization identity is
\[
                       \psi\lambda_M\phi=D\lambda_C^{-1}.
\]
It gives the first numerical equality; applying norm and cancelling
$D$ in the torsion-free Neron--Severi group gives the second.
Isogeny invariance of $2$-rank gives $f_B=1$.

Conversely, an etale isogeny induces an isomorphism on the
local-local $2$-divisible factors.  The local-local elementary
divisors of $\phi^*M$ are therefore both $b(M)$.  Since $\Theta_B$
is principal, those of $D\Theta_B$ are both $v_2(D)$.
The equality of their polarizations implies $b(M)=v_2(D)$,
and hence $a(M)=0$.
\end{proof}

\begin{lemma}[The even-degree mixed Prym bound without separability]
\label{lem:DE-A3-mixed-all-even-Prym}
\status{Proved; root and monodromy full independent proof/source/code audits passed.}
Let $A,M,D$ be as in
Lemma~\ref{lem:DE-A3-mixed-polarization-blocks}, with no restriction
on either index.  Every smooth curve $Y\in|M^m|$ with $m$ even
satisfies
\[
                              g(Y)-f_Y\ge2.
\]
\end{lemma}
\begin{proof}
The proof of Lemma~\ref{lem:DE-A3-general-polarization-Prym}
uses oddness of $D$ only to construct the integral line bundle $N$.
That construction is now supplied for every $M$ by
Lemma~\ref{lem:DE-A3-mixed-polarization-blocks}.
We give the remaining argument explicitly.
Ample vanishing and Serre duality give
$H^1(A,P\otimes M^{-m})=0$ for every $P\in\Pic^0(A)$.
For nontrivial $P$, also $H^0(A,P)=0$, so restriction to $Y$
cannot trivialize $P$.  The same vanishing for $P=\OO_A$ makes
the tangent map on Picard varieties injective.  Thus restriction
is an immersion $j:A^\vee\hookrightarrow J_Y$.
Its polarized dual $q:J_Y\to A$ has smooth connected kernel $P_Y$:
quotient by the reduced identity kernel and use isogeny duality
to exclude a remaining finite quotient, exactly as in the cited lemma.

The integral polarization identity is
\[
                 qj=j^*\lambda_{J_Y}=mD\lambda_M^{-1}=m\lambda_N.
\]
Here the first equality identifies the polarization with its
homomorphism $A^\vee\to A$; the middle identity is
\cite[Section~3, equations~(3.2)--(3.3), pp.~346--347]
{Debarre1994Degrees}.  It holds in all characteristics.
Since $m$ is even, the actual finite group scheme $A^\vee[2]$
embeds in $P_Y[2]$.
Twist this finite-group immersion by Frobenius and apply functorial
Verschiebung.  The connected part of
$\ker(V_{A^\vee}:(A^\vee)^{(2)}\to A^\vee)$ has order two and injects
into the connected part of
$\ker(V_{P_Y}:P_Y^{(2)}\to P_Y)$.
Hence $P_Y$ has $2$-rank at most $\dim P_Y-1$.
The quotient $J_Y/P_Y\simeq A$ has defect one; adding defects,
using isogeny invariance and a complementary abelian subvariety,
gives $g(Y)-f_Y\ge2$.
\end{proof}

\begin{theorem}[All degrees for mixed surfaces of ordinary index zero]
\label{thm:DE-A3-mixed-nonproduct-ranks}
\status{Proved; root and monodromy full independent proof/source/code audits passed.}
Let $A$ be an abelian surface of $2$-rank one over a perfect field
of characteristic two, and let $M$ be ample.  Compute its ordinary
index after passage to an algebraic closure, and assume $a(M)=0$.
Write $D=\chi(M)$ and $L=M^3$.  Then $L$ is very ample.
For every $d\ge1$, the geometric generic relative stable rank
in the complete system $|L^d|$ is
\[
 r_d=
 \begin{cases}
       9Dd^2-1,&d\ \text{odd},\\
       9Dd^2-2,&d\ \text{even}.
 \end{cases}
\]
The relative coherent dimension is $9Dd^2-1$, and the generic
smooth curve has genus $9Dd^2+1$.  For the embedding defined by $L$,
the same ranks hold for ambient hypersurfaces whenever
$d\ge\max\{1,\operatorname{reg}(\mathcal I_A)-1\}$.
\end{theorem}
\begin{proof}
Work over an algebraic closure.  The etale source from
Lemma~\ref{lem:DE-A3-mixed-index-zero-source}, and the norm
point-and-tangent argument in
Theorem~\ref{thm:DE-A3-ordinary-all-ranks}, show that $M^3$ is
very ample.  That argument works for any finite etale degree:
one chooses a section with the desired value or first derivative
at one chosen lift, avoids its finitely many other value hyperplanes,
and takes its norm.  Numerical equivalence is converted to a
translation because the ample polarization is surjective.

For odd $d$, the ordinary index of $M^d$ is still zero.
Apply Lemma~\ref{lem:DE-A3-mixed-index-zero-source} to $M^d$.
Its principal source $B_d$ has $2$-rank one, and the etale norm
degree is $Dd^2$.  Proposition~\ref{prop:DE-A3-ppav-degree-one}
gives a nonempty open of smooth genus-ten curves of $2$-rank nine
in $|3\Theta_{B_d}|$.
A general such curve has integral nodal image under the isogeny.
The full cubic nodal-image proof in
Theorem~\ref{thm:DE-A3-ordinary-all-ranks} applies here:
its two-first-jet and three-value incidence arguments concern
the principal source and distinct kernel points, and do not use
ordinarity of the source or invertibility of the kernel order.
Thus the image has $9(Dd^2-1)$ nodes and normalization of defect one.
An arbitrary target translation puts its line bundle in the exact
system $|M^{3d}|$, since its norm class is $3dM$.
This uses surjectivity of the polarization, not an inverse on
two-torsion points.

The normalization sequence adds a Frobenius-bijective graph space.
The nodal image therefore has full defect one and the maximal
possible relative stable rank $9Dd^2-1$.
Ample vanishing gives constant full coherent dimension
$9Dd^2+1$ and relative dimension $9Dd^2-1$ over the coefficient
space.  A sufficiently long Frobenius iterate detects the stable
rank.  Its maximal-rank open meets the nonempty smooth-member open,
proving the odd formula.

For even $d$, choose sharp smooth members of $|L|$ and
$|L^{d-1}|$ meeting transversely.  Their connected nodal union
has full defect two, because each normalization component has
defect one and the graph contribution is Frobenius-bijective.
Openness gives the lower bound $r_d\ge9Dd^2-2$.
Lemma~\ref{lem:DE-A3-mixed-all-even-Prym}, applied with $m=3d$,
gives the matching upper bound for every smooth member.
Finally the stated regularity bound makes restriction from ambient
hypersurfaces surjective, with geometrically integral parameter
fibres, and hence preserves the generic rank.
\end{proof}

\begin{remark}[What the ordinary index adds]
\label{rem:DE-A3-mixed-nonproduct-scope}
\status{Proved; root and monodromy full independent proof/source/code audits passed.}
The condition in Theorem~\ref{thm:DE-A3-mixed-nonproduct-ranks}
is intrinsic to the polarized surface and does not require an
elliptic product or an odd polarization degree.
For a product of an ordinary and a supersingular elliptic curve,
with product polarization of degrees $(u,v)$ in that order,
one has $a(M)=v_2(u)$ and $b(M)=v_2(v)$.
Thus this theorem includes the odd-ordinary-factor product cases,
while also applying to any nonproduct surface satisfying the
same integral condition.
Such even-degree classes can be constructed directly: start with
any principally polarized surface of $2$-rank one and quotient by
an etale subgroup of order $2^b$ in its rank-one etale
$2$-divisible factor.  The norm of the principal bundle is ample,
has Euler characteristic $2^b$, and its pullback has numerical
class $2^b\Theta$.  The converse in
Lemma~\ref{lem:DE-A3-mixed-index-zero-source} therefore gives
ordinary index zero.  If the chosen source is geometrically simple,
its isogenous target remains so.
The class $a(M)>0$ is not settled by this norm construction:
Lemma~\ref{lem:DE-A3-mixed-index-zero-source} shows precisely why
an etale principal source of the prescribed degree cannot exist
there.  No formula for arbitrary ambient varieties is asserted.

The retained certificate
\path{agents/obstructions_crystals/code/mixed_nonproduct_rank_checks.py}
proves the block inverse and determinant identities symbolically,
including an arbitrary unit multiplier.  It also checks $289$
index pairs, $1024$ odd norm identities and $1024$ even union identities.
The author run passed in $0.04614707099972293$ seconds, with its
complete output retained alongside the program.  The existence and
group-scheme assertions follow from the proofs above, not from the
finite arithmetic checks.
Root independently read the entire proof, original sources and program,
and reran the certificate in $0.049978896997345146$ seconds.
Monodromy independently completed the same proof/source/program audit
and reran it in $0.03037099799985299$ seconds.  Their outputs are retained
separately; the author output was preserved.
\end{remark}
